\begin{abstract}
Autoresearch systems iterate an LLM-driven propose--train--evaluate loop to
improve a target task, but their search mechanisms are fixed at design time.
Human researchers introduce improvements such as multi-batch parallelism or
persistent memory; the system itself cannot.
We present \methodname{}, a two-level framework in which an outer loop
autonomously discovers and injects new search mechanisms into the inner loop.
Level~1 runs the standard autoresearch loop; Level~1.5 analyzes the search
trace and adjusts active parameters; Level~2 conducts a 4-round research
session (Explore, Critique, Specify, Generate) and writes Python code that
structurally modifies the inner loop at runtime via \texttt{importlib}.
We evaluate four configurations on Karpathy's GPT pretraining benchmark
(\valbpb{}) using the same DeepSeek model at every level.
The pure inner loop (Group~A) achieves $-0.009 \pm 0.002$; adding Level~1.5 alone (Group~B) yields no reliable additional gain ($-0.006 \pm 0.006$).
Adding Level~2 (Group~C) achieves $-0.045 \pm 0.030$, a 5$\times$ improvement
over Group~A in two of three repeats (with high variance: $\pm 0.030$).
Group~D (Level~1 + Level~2, without Level~1.5) achieves $-0.034 \pm 0.031$, confirming that Level~2 alone drives the improvement.
Across three independent repeats, Level~2 autonomously generates mechanisms
from three distinct domains---combinatorial optimization, multi-armed bandits,
and design of experiments (a fourth domain, Bayesian optimization, was attempted
but reverted due to a missing dependency)---without human specification.
These results validate the principle that \emph{autoresearch can research
itself}: structural modification of the search loop, not merely parameter
adjustment, is the decisive lever for large performance gains.
\end{abstract}
